\documentclass[11pt]{article}
\usepackage[margin=1in]{geometry}
\usepackage{amsmath}
\usepackage{booktabs}
\usepackage{microtype}
\usepackage{hyperref}
\usepackage{natbib}
\usepackage{enumitem}
\usepackage{xurl}

\title{LAM-JEPA on ARC-Challenge: A Reproducible Falsification-First Evaluation}
\author{[AUTHORS PENDING OWNER APPROVAL]}
\date{}

\begin{document}
\maketitle

\begin{abstract}
We evaluate the project-named LAM-JEPA system on ARC-Challenge using a frozen external-benchmark pipeline, a gradient-active-parameter-matched supervised comparator, mechanism ablations, a shuffled-label control, and a pinned pretrained comparator. Source inspection materially changes how the experiment should be described: the frozen ARC path is a small hashed-token, mean-pooled embedding model with vector quantization, learned sparse memory, a one-step latent-action rollout, and same-input exponential-moving-average (EMA) target alignment; it is not a Transformer encoder and it does not instantiate the canonical I-JEPA context-to-distinct-target prediction task. Across five frozen validation seeds, the full model achieved accuracy $0.2549152542 \pm 0.0129968064$, while the matched supervised model achieved $0.2664406780 \pm 0.0154600058$, for a paired LAM-minus-matched difference of $-0.0115254237 \pm 0.0140994131$. Planner and target-path ablations likewise failed their preregistered contribution criteria. A later trainability repair passed its bounded train-only gate but repaired validation remained negative/inconclusive. We preserve these outcomes, keep the confirmatory test locked, and report the experiment as a reproducible falsification case study rather than an architecture-superiority result.
\end{abstract}

\section{Introduction}
Representation-learning projects are vulnerable to a common failure mode: an interesting mechanism is treated as validated because the code runs, a favorable development slice exists, or a later repair improves trainability. This work asks a narrower question: under a frozen ARC-Challenge validation protocol, does the tested LAM-JEPA configuration outperform an appropriate capacity-matched supervised baseline, and do its planner and target-path mechanisms contribute measurably under preregistered criteria?

The answer on the current evidence is no. The purpose of this manuscript is therefore not to report a superiority result. It documents a reproducible negative/inconclusive evaluation with enough protocol detail, controls, raw evidence, source inspection, and claim boundaries to make the falsification useful.

Our contributions are: (1) a reproducible ARC-Challenge evaluation path with retained eligibility and exclusion evidence; (2) a gradient-active-parameter-matched supervised comparison; (3) five-seed planner/target ablations and a deterministic shuffled-label control; (4) a bounded pinned pretrained-comparator path; (5) a documented trainability repair whose later validation did not rescue the original scientific claim; (6) an explicit stop rule forbidding use of the locked ARC confirmatory test to rescue the failed hypothesis; and (7) a source-level audit that constrains architecture and novelty wording to what was actually tested.

No claim of ARC superiority, planner benefit, target-path benefit, quantization benefit, Transformer reasoning capability, general JEPA failure, general benchmark superiority, or research completeness is made.

\section{Related Work and Novelty Boundary}
I-JEPA predicts representations of distinct target image blocks from a context block and updates a target encoder by exponential moving average \citep{assran2023ijepa}. The frozen LAM ARC path differs materially: its online and target encoders receive the same serialized ARC input, and the ARC loss aligns the quantized online latent to that same-input EMA target. We therefore use conservative \emph{target-encoder alignment} wording and do not claim to instantiate the canonical I-JEPA context-to-distinct-target task.

Vector-quantized discrete latent learning predates this project \citep{vandenOord2017vqvae}. Latent Action Pretraining learns discrete latent actions from video \citep{ye2024lapa}, while recent latent-action world-model work studies constrained latent actions and planning in action-free video settings \citep{garrido2026latentaction}. FF-JEPA introduces an action-free latent planner alongside forward dynamics for long-horizon planning \citep{masip2026ffjepa}. These directions make generic novelty claims about discrete latent actions or latent planners inappropriate here.

ARC was introduced as a multiple-choice science-question benchmark intended to stress reasoning beyond earlier QA datasets \citep{clark2018arc}. Our bounded pretrained comparison uses a pinned DeBERTa-v3-xsmall checkpoint from the DeBERTaV3 family \citep{he2023debertav3}. Reproducibility programs and checklists are established in machine learning \citep{pineau2021reproducibility}; our contribution is the use of frozen protocols, adverse-result retention, claim ledgers, artifact lineage, and explicit stop rules for this particular negative result.

\section{Hypotheses and Falsification Rule}
Let $A_{\mathrm{full}}$ denote validation accuracy of the frozen full configuration, $A_{\mathrm{match}}$ the gradient-active-parameter-matched supervised comparator, $A_{\mathrm{np}}$ the \texttt{no\_planner} ablation, and $A_{\mathrm{nt}}$ the \texttt{no\_target} ablation. The frozen line tested three primary claims:
\begin{enumerate}[leftmargin=*]
    \item \textbf{H1---superiority:} $A_{\mathrm{full}}$ should outperform $A_{\mathrm{match}}$ under the frozen ARC validation protocol.
    \item \textbf{H2---planner contribution:} removing the planner should produce a sufficiently adverse paired effect under the preregistered mechanism criterion.
    \item \textbf{H3---target-path contribution:} replacing the EMA target path with the frozen \texttt{no\_target} control should produce a sufficiently adverse paired effect under the preregistered mechanism criterion.
\end{enumerate}
The deterministic shuffled-label control was required to remain below a frozen accuracy ceiling of $0.35$. It is a validity control, not evidence for H1--H3.

\paragraph{Falsification rule.} If frozen validation fails the superiority/mechanism criteria, the locked ARC confirmatory test is not opened to rescue the line. Architecture repairs or successor mechanisms become new versioned hypotheses.

\section{Method}
\subsection{ARC serialization and encoder}
Each ARC example is formatted as a question followed by indexed answer choices. The frozen \texttt{text\_to\_tokens} path lowercases the string, splits on whitespace, hashes each token deterministically with BLAKE2b, maps the hash modulo a vocabulary of 256 IDs, and pads or truncates to 96 positions. The ARC adapter supplies a zero-valued numeric vector for every example, so the numeric branch provides no example-varying ARC information.

The frozen token encoder is not a Transformer. It uses a learned token embedding (vocabulary 256, embedding dimension 32), learned positional vectors, an identity sequence encoder, layer normalization, and mean pooling. Token and numeric representations are fused by an MLP, normalized, and projected to a 32-dimensional latent $z$.

\subsection{Quantizer, memory, planner, and target path}
The latent passes through a 32-code vector quantizer with nearest-code assignment, EMA codebook statistics, commitment/codebook MSE, and a straight-through estimator, producing $z_q$. A learned sparse memory retrieves a top-$k$ weighted value correction. The planner contains an eight-action policy, learned action embeddings, and residual transition networks. The frozen retained protocol uses one model step, followed by a four-choice answer head.

The model maintains a target encoder/projector initialized from the online encoder/projector and updated by EMA with momentum $0.996$. In the frozen ARC forward pass, the target receives the same tokens and numeric input as the online encoder. The \texttt{no\_target} ablation replaces the EMA target representation with detached online $z$ rather than removing alignment entirely.

\subsection{Objective and matched comparator}
For an ARC minibatch, the reported benchmark uses
\[
L = L_{\mathrm{CE}} + 0.5L_{\mathrm{align}} + 0.25L_{\mathrm{quant}} + 0.25L_{\mathrm{traj}},
\]
where $L_{\mathrm{CE}}$ is four-choice cross entropy, $L_{\mathrm{align}}$ is cosine alignment between $z_q$ and detached EMA target $z_{\mathrm{target}}$, $L_{\mathrm{quant}}$ is the quantizer loss, and $L_{\mathrm{traj}}$ is mean squared error from planner rollout states to detached $z_q$.

The matched supervised comparator uses the same encoder/projector family and a four-choice classifier without target, quantizer, memory, or planner machinery. Gradient-active parameter counts are 86,372 for LAM-JEPA and 86,644 for the matched supervised model, a ratio of 1.0031491687.

\section{Experimental Setup}
The frozen eligibility rule retained 1,117 of 1,119 ARC-Challenge training rows and 295 of 299 validation rows. Excluded rows are retained as evidence. The locked ARC test was not downloaded or evaluated for this failed hypothesis line.

The full-controls validation uses seeds $\{1,2,3,4,5\}$, 20 epochs, batch size 32, learning rate $3\times10^{-4}$, one model step, all eligible train and validation rows, and CPU execution. The pinned pretrained comparator uses \texttt{microsoft/deberta-v3-xsmall} at immutable revision \texttt{14809e4f1fe1895fcba8b258271a940c6ca45ec4}. Means and sample standard deviations are reported across seeds; mechanism effects are paired by seed and summarized with retained bootstrap 95\% intervals.

\section{Results}
\subsection{Capacity-matched comparison}
\begin{table}[h]
\centering
\caption{Frozen ARC-Challenge validation result.}
\begin{tabular}{lr}
\toprule
System & Validation accuracy \\
\midrule
LAM-JEPA & $0.2549152542 \pm 0.0129968064$ \\
Matched supervised & $0.2664406780 \pm 0.0154600058$ \\
Paired LAM $-$ matched & $-0.0115254237 \pm 0.0140994131$ \\
\bottomrule
\end{tabular}
\end{table}
The frozen validation does not support superiority over the capacity-matched supervised baseline.

\subsection{Mechanism ablations}
\begin{table}[h]
\centering
\caption{Frozen mechanism and validity controls.}
\begin{tabular}{lr}
\toprule
Configuration & Validation accuracy \\
\midrule
Full LAM-JEPA & $0.2549152542 \pm 0.0129968064$ \\
\texttt{no\_planner} & $0.2501694915 \pm 0.0129968064$ \\
\texttt{no\_target} & $0.2616949153 \pm 0.0203954020$ \\
Shuffled-label control & $0.2630508475 \pm 0.0145011862$ \\
\bottomrule
\end{tabular}
\end{table}
Full minus \texttt{no\_planner} was $+0.0047457627$ with retained bootstrap 95\% CI $[0.0, 0.0142372881]$. Full minus \texttt{no\_target} was $-0.0067796610$ with interval $[-0.0135593220,0.0]$. Neither required mechanism criterion is supported. The shuffled-label control remained below the frozen 0.35 ceiling but is not presented as favorable evidence for the architecture.

\subsection{Bounded pretrained comparison and repaired validation}
On the recorded bounded development comparison, LAM-JEPA achieved 0.15625 and DeBERTa-v3-xsmall achieved 0.21875, a paired difference of $-0.0625$. This is characterization evidence rather than a broad inferiority claim.

The independent recomputation verdict for repaired ARC-v5 validation is \texttt{VALID\_NEGATIVE\_OR\_INCONCLUSIVE\_VALIDATION}. The trainability repair passed its bounded train-only gate, but repaired validation did not establish generalization or quantization benefit. Repaired-minus-legacy paired mean was $+0.0040678024$ with bootstrap 95\% CI $[-0.0135593116,0.0216949165]$; repaired-minus-no-quantizer was $+0.0033898324$ with bootstrap 95\% CI $[-0.0027118593,0.0094915211]$.

\section{Robustness and Failure Analysis}
The frozen controls were independently rerun from scientific source SHA \texttt{760aa7f9a73a177d5ff4ba7eb470f7e68ace63cb}. Retained successful attempts include GitHub Actions run 31203337502 attempt 2 (artifact 9149336081; SHA-256 \texttt{c45710b5dae6a767ccb6bab7f6e3d8e9578752d8cf9b79fd82a65ae824dded1b}) and attempt 3 (artifact 9162165932; SHA-256 \texttt{caa898f1ff046a337db9b5ddbffe1b332943a732868e2fd809abeda8ee89c30b}). Eight of ten retained files were byte-identical. Raw result/input JSON showed low-order floating-point drift, with maximum observed numeric difference approximately $5.9186\times10^{-4}$, while aggregate accuracies, mechanism effects, verifier reports, and the scientific conclusion were unchanged.

A source audit identifies limitations that plausibly contribute to the adverse result but must not be used as rescue arguments: ARC text is represented by deterministic whitespace-token hashes into only 256 token IDs; the encoder mean-pools embeddings without attention or contextual sequence processing; and the numeric branch is constant across ARC examples. The EMA target also receives the same serialized input as the online encoder, so the alignment objective does not directly train prediction of a held-out future or masked target.

A later train-only investigation localized a failure in the quantized latent path. The opt-in repair \texttt{arc-v5-stable-ema-residual-0.03125} passed its bounded trainability gate and was independently reproduced before repaired validation, which remained negative/inconclusive. Separately, a pre-fix seed-order defect was found because model initialization occurred before the requested seed was applied; PR \#61 / SHA \texttt{b72a97a99769b278eb8ec75bc5eab62dc9599f29} repaired reproducibility plumbing without changing the frozen scientific protocol.

\section{Limitations}
The conclusion is specific to the frozen ARC line and tested small configurations. The frozen token encoder is not a Transformer or modern pretrained language encoder. The ARC target path uses same-input EMA alignment rather than distinct context/target prediction. Five seeds quantify only a narrow protocol. The bounded DeBERTa comparison is not a full matched pretrained-comparator study. The repaired quantization path changes trainability but does not validate the original mechanism claim. Exact physical CPU model metadata is not retained in the paper-facing evidence. Independent outside expert review is still absent. License, authorship, and final citation metadata remain owner-controlled release gates. Finally, the locked confirmatory test cannot ethically be used as a rescue set after the validation hypothesis failed.

\section{Reproducibility and Provenance}
Every quantitative manuscript statement is required to trace through the chain
\[
\text{claim} \rightarrow \text{table/figure} \rightarrow \text{processed metric} \rightarrow \text{raw retained artifact} \rightarrow \text{frozen config/protocol} \rightarrow \text{source commit}.
\]
The repository's \texttt{MANUSCRIPT\_PROVENANCE.md} and \texttt{REPRODUCE.md} contain the current provenance edges, exact commands, artifact identifiers, and frozen stop rules. The defensible reproducibility statement is that independent project-controlled reruns reproduce the aggregate scientific conclusion and verifier outputs; low-level probabilities and checkpoint bytes are not claimed bitwise identical across all runners. Independent external reproduction remains pending.

\section{Discussion}
The most useful outcome is that the evidence pipeline prevented an executable research prototype from being mislabeled as a successful research result. Capacity matching removed one easy confound, while mechanism ablations prevented the full-model score from being attributed automatically to the planner or EMA target path. The trainability repair demonstrates why engineering recovery and scientific validation must remain separate: making optimization behave better did not produce the expected validation advantage.

The source audit also changes how the system should be described. In the frozen ARC path, ``LAM-JEPA'' is a project identifier, not proof that a canonical JEPA prediction problem or Transformer reasoning model was evaluated. This distinction makes the negative report more useful because readers can see exactly what failed and what remains untested.

\section{Conclusion}
Under the frozen ARC-Challenge validation protocol, the tested LAM-JEPA configuration did not outperform the capacity-matched supervised baseline and did not validate its planner or EMA-target contribution criteria. A later trainability repair also failed to produce a positive repaired-validation verdict. The result is intentionally narrow: it falsifies the current frozen ARC claims, not JEPA, vector quantization, latent planning, or representation learning in general. Adverse results are retained as first-class artifacts, and the confirmatory test remains locked for this failed line.

\paragraph{Release status.} The scientific manuscript and internal reproducibility package are complete for the frozen evidence. Public release still requires owner-approved authorship/license metadata and, if claimed, genuinely independent external reproduction. This source intentionally does not invent those fields.

\bibliographystyle{plainnat}
\bibliography{references}

\end{document}
